\documentclass[11pt,a4paper]{article}
\usepackage[T1]{fontenc}
\usepackage{lmodern}
\usepackage{microtype}
\usepackage{amsmath,amssymb,amsthm,mathtools,mathrsfs}
\usepackage[margin=28mm]{geometry}
\usepackage{aliascnt}
\usepackage{needspace}
\usepackage[hidelinks]{hyperref}
\usepackage[nameinlink,noabbrev]{cleveref}
\hypersetup{
 pdftitle={Residual conductivity at isotropic conical band crossings},
 pdfsubject={A proof of the residual conductivity theorem for a periodic semimetal},
 pdfkeywords={graphene, conductivity, conical crossings, linear response}
}

\newcommand{\R}{\mathbb R}
\newcommand{\C}{\mathbb C}
\newcommand{\Z}{\mathbb Z}
\newcommand{\B}{\mathcal B}
\newcommand{\lat}{\mathcal R}
\newcommand{\T}{\mathbb T_{\!*}^{2}}
\newcommand{\HH}{\mathscr H}
\newcommand{\DD}{\mathscr D}
\newcommand{\one}{\mathbf 1}
\newcommand{\trv}{\underline{\operatorname{Tr}}}
\newcommand{\ip}[2]{\langle #1,#2\rangle}
\newcommand{\norm}[1]{\lVert #1\rVert}
\newcommand{\abs}[1]{\lvert #1\rvert}
\newcommand{\cB}{\frac{1}{(2\pi)^2}}
\DeclareMathOperator{\Tr}{Tr}
\DeclareMathOperator{\tr}{tr}
\DeclareMathOperator{\Ran}{Ran}
\DeclareMathOperator{\Ree}{Re}
\DeclareMathOperator{\Imm}{Im}

\newtheorem{theorem}{Theorem}[section]
\newaliascnt{proposition}{theorem}
\newtheorem{proposition}[proposition]{Proposition}
\aliascntresetthe{proposition}
\newaliascnt{lemma}{theorem}
\newtheorem{lemma}[lemma]{Lemma}
\aliascntresetthe{lemma}
\crefname{theorem}{Theorem}{Theorems}
\crefname{proposition}{Proposition}{Propositions}
\crefname{lemma}{Lemma}{Lemmas}
\crefname{section}{Section}{Sections}
\crefname{appendix}{Appendix}{Appendices}
\Crefname{proposition}{Proposition}{Propositions}
\Crefname{lemma}{Lemma}{Lemmas}
\numberwithin{equation}{section}
\setlength{\emergencystretch}{2em}

\title{Residual conductivity at isotropic conical band crossings}
\author{}
\date{}

\begin{document}
\maketitle
\vspace{-1.5em}
\begin{abstract}
We consider spinless, noninteracting electrons governed by a real periodic
Schr\"odinger operator in two dimensions. Suppose that the Fermi energy
meets the spectrum at finitely many isolated, isotropic conical crossings
of two adjacent bands. For a constant electric field switched on at time
zero, we prove that the Ces\`aro mean of the linear-response conductivity
is $1/16$ times the identity per crossing, in atomic units. The proof
starts from the full Schr\"odinger evolution. An integrated band-energy
identity cancels the term linear in time, time reversal cancels the
antisymmetric response, and a local two-dimensional spectral reduction
determines the remaining contribution. The conical eigenvalue assumptions
alone imply the required current matrix elements. We also justify the
Stark dynamics, the trace per unit area, and differentiation of the
unbounded current for potentials in $H^1_{\mathrm{per}}$.
\end{abstract}

\section{The model and the theorem}

At an isolated conical crossing, a two-dimensional semimetal has no Fermi
curve, but its response to a weak electric field can have a nonzero
long-time average. We prove this statement for a periodic Schr\"odinger
operator under explicit assumptions on its bands. The order of limits is
part of the result: we first take the derivative at zero field at each
fixed time, and then average that derivative over an increasing time
interval.

Let $\lat\subset\R^2$ be a lattice, let $\Omega$ be a fundamental cell,
and define the reciprocal lattice by
\[
 \lat^*=\{K\in\R^2:K\cdot R\in2\pi\Z\text{ for all }R\in\lat\}.
\]
We write $\B$ for a reciprocal fundamental cell and
$\T=\R^2/\lat^*$ for the Brillouin torus. In particular,
\begin{equation}\label{eq:volumes}
 |\Omega|\,|\B|=(2\pi)^2.
\end{equation}
All scalar products are linear in the second argument. Electrons are
spinless, and we use atomic units with the signs of the field and current
specified below.

Let $V\in H^1_{\mathrm{per}}(\Omega;\R)$ and put $p=-i\nabla$. The
Hamiltonian and its periodic Bloch fibers are
\begin{equation}\label{eq:hamiltonians}
 H=\tfrac12p^2+V\quad\text{on }L^2(\R^2),
 \qquad
 H_k=\tfrac12(p+k)^2+V\quad\text{on }\HH:=L^2_{\mathrm{per}}(\Omega).
\end{equation}
Their domains are $H^2(\R^2)$ and
$\DD:=H^2_{\mathrm{per}}(\Omega)$, respectively. Each fiber has compact
resolvent. Denote its ordered eigenvalues, counted with multiplicity, by
$\lambda_1(k)\leq\lambda_2(k)\leq\cdots$, and set
$\lambda_0=-\infty$.

Fix a Fermi energy $\mu_{\mathrm F}\in\R$. We assume that there are an
integer $N\geq1$ and distinct points $k_1,\ldots,k_M\in\T$ such that
\begin{equation}\label{eq:fermi-gap}
 \lambda_N(k)<\mu_{\mathrm F}<\lambda_{N+1}(k)
 \qquad\text{for }k\in\T\setminus\{k_1,\ldots,k_M\},
\end{equation}
whereas at each crossing,
\begin{equation}\label{eq:isolated-double}
 \lambda_{N-1}(k_i)<\lambda_N(k_i)=\mu_{\mathrm F}
 =\lambda_{N+1}(k_i)<\lambda_{N+2}(k_i).
\end{equation}
Finally, for some $v_i>0$,
\begin{equation}\label{eq:cones}
 \begin{split}
 \lambda_N(k_i+q)&=\mu_{\mathrm F}-v_i|q|+O(|q|^2),\\
 \lambda_{N+1}(k_i+q)&=\mu_{\mathrm F}+v_i|q|+O(|q|^2)
 \qquad(q\longrightarrow0).
 \end{split}
\end{equation}
These are local assertions in Euclidean coordinates on $\T$. The
remainder bounds are uniform in the direction of $q$. No assumption is
made about crossings among occupied bands or among unoccupied bands.

Let $e_\beta,e_\alpha\in\R^2$ specify the field and current directions;
they need not be unit vectors. For $\varepsilon\geq0$, define
\begin{equation}\label{eq:stark-density}
 \begin{gathered}
 \gamma_0=\one_{(-\infty,\mu_{\mathrm F}]}(H),\qquad
 H_\varepsilon=H-\varepsilon x\cdot e_\beta,\qquad
 J_\alpha=p\cdot e_\alpha,\\
 \gamma_\varepsilon(t)
 =e^{-itH_\varepsilon}\gamma_0e^{itH_\varepsilon},\qquad
 j_\varepsilon(t)=\trv(J_\alpha\gamma_\varepsilon(t)).
 \end{gathered}
\end{equation}
Here $H_\varepsilon$ is the self-adjoint closure of its restriction to
$C_c^\infty(\R^2)$, and the trace per unit area of a periodic, locally
trace-class operator is
\begin{equation}\label{eq:trace-definition}
 \trv(A)=\frac1{|\Omega|}
 \Tr\bigl(\one_\Omega A\one_\Omega\bigr).
\end{equation}
The existence of this closure and of the current trace in
\eqref{eq:stark-density} is proved in \cref{sec:appendix}.

\begin{theorem}\label{thm:main}
Under \eqref{eq:fermi-gap}--\eqref{eq:cones}, the current
$j_\varepsilon(t)$ is well defined for every $\varepsilon,t\geq0$, and
$j_0(t)=0$. For every fixed $t\geq0$, the limit
\[
 \sigma_{\alpha\beta}(t)
 :=\lim_{\varepsilon\downarrow0}\frac{j_\varepsilon(t)}{\varepsilon}
\]
exists. Moreover,
\begin{equation}\label{eq:main-result}
 \lim_{T\longrightarrow\infty}\frac1T\int_0^T
 \sigma_{\alpha\beta}(t)\,dt
 =\frac{M}{16}\,e_\alpha\cdot e_\beta.
\end{equation}
\end{theorem}

Each distinct crossing on the reciprocal torus is counted once. The
motivating case is the two-crossing band geometry of spinless graphene;
it gives $1/8$ times the identity. The
statement concerns the iterated limit in \eqref{eq:main-result}; it does
not require pointwise convergence of $\sigma_{\alpha\beta}(t)$ as
$t\to\infty$.

The proof has three parts. We first derive the fixed-time response formula
for the full Hamiltonian. We then use the total energy of
the occupied bands to cancel the apparent linear growth in time. Finally,
we reduce the contribution of each crossing to an elementary integral.
The standard operator results used in the argument are drawn from
Reed--Simon \cite{RS1,RS2,RS4}; the analytic details needed here are
included in the appendix.

\section{Bloch dynamics and linear response}\label{sec:response}

We use the Bloch transform
\begin{equation}\label{eq:bloch-transform}
 (\mathcal U f)(k,x)
 =\sum_{R\in\lat}e^{-ik\cdot(x+R)}f(x+R),
 \qquad x\in\Omega,
\end{equation}
with direct-integral measure $dk/|\B|$ and the ordinary $L^2(\Omega,dx)$
norm on each fiber. Lattice Parseval gives its unitarity. In this
representation, $H$ and $J_\alpha$ have fibers $H_k$ and $v_\alpha(k)$,
respectively; these identities follow on a core and then by closure
\cite[Theorem XIII.97]{RS4}. The trace per unit area is
\begin{equation}\label{eq:bloch-trace}
 \trv(A)=\cB\int_\B\Tr_{\HH}A(k)\,dk
\end{equation}
for the operators used below. The factor follows from
\eqref{eq:volumes}; a trace-class justification is given in
\cref{lem:local-trace}.

Set
\begin{equation}\label{eq:projector-velocity}
 P(k)=\one_{(-\infty,\mu_{\mathrm F}]}(H_k),\qquad
 v_\eta(k)=(p+k)\cdot e_\eta=\partial_\eta H_k,
 \quad \eta\in\{\alpha,\beta\},
\end{equation}
where $\partial_\eta=e_\eta\cdot\nabla_k$. The projection $P(k)$ has
rank $N$ almost everywhere. Its rank is $N+1$ at a crossing, but changes
on this finite set do not affect any direct integral or current. All
band formulas below are understood away from these points.

For $K\in\lat^*$, the unitary $\tau_Ku(x)=e^{-iK\cdot x}u(x)$ satisfies
\[
 H_{k+K}=\tau_KH_k\tau_K^*,\qquad
 v_\eta(k+K)=\tau_Kv_\eta(k)\tau_K^*.
\]
Consequently, band energies and the traces below are functions on $\T$.

Choose an orthonormal eigenbasis $u_n(k)$, measurable in $k$, and define
for $n\leq N<m$
\begin{equation}\label{eq:gap-matrix}
 \begin{split}
 \Delta_{mn}(k)&=\lambda_m(k)-\lambda_n(k),\\
 z_{nm}(k)&=
 \ip{u_n(k)}{v_\alpha(k)u_m(k)}
 \ip{u_m(k)}{v_\beta(k)u_n(k)}.
 \end{split}
\end{equation}
The response and sum-rule formulas below are independent of the choice
of basis within a degenerate eigenspace; the absolute-value estimates
hold for every such choice. Existence of a measurable
eigenbasis is part of the Bloch spectral theory in
\cite[Theorem XIII.98]{RS4}.

\subsection{Bounds for the band sums}

The uniform graph estimates in \cref{lem:operators} and the boundedness
of the occupied eigenvalues imply
\begin{equation}\label{eq:occupied-bound}
 \sup_{k\in\B}\sum_{n=1}^N
 \norm{v_\eta(k)u_n(k)}^2<\infty.
\end{equation}
Throughout, a supremum involving $P(k)$ may be read as an essential
supremum. By Cauchy--Schwarz and Parseval,
\begin{equation}\label{eq:sum-bound}
 \sum_{n\leq N<m}|z_{nm}(k)|
 \leq\sum_{n=1}^N
 \norm{v_\alpha(k)u_n(k)}\,
 \norm{v_\beta(k)u_n(k)}\leq C.
\end{equation}
In particular, this estimate controls the infinitely many unoccupied
bands without an energy cutoff.

Choose disjoint small disks $D_i$ about the crossings on $\T$. On the
complement of these disks, all occupied-to-unoccupied gaps have a common
positive lower bound. Inside $D_i$, the only gap that can become small
is
\begin{equation}\label{eq:small-gap}
 \Delta_i(q):=\Delta_{N+1,N}(k_i+q)
 =2v_i|q|+O(|q|^2),\qquad
 c|q|\leq\Delta_i(q)\leq C|q|.
\end{equation}
Every other occupied-to-unoccupied gap in these disks is bounded below
uniformly. These assertions follow from continuity of the bands,
compactness of the torus, and \eqref{eq:isolated-double}. Hence
\begin{equation}\label{eq:inverse-gap-integrable}
 k\longmapsto
 \sum_{n\leq N<m}\frac{|z_{nm}(k)|}{\Delta_{mn}(k)}
 \quad\text{belongs to }L^1(\B).
\end{equation}
The possible singularity is $C/|k-k_i|$, which is integrable in two
dimensions.

\subsection{The fixed-time response}

For each fiber, let $U_{\varepsilon,k}(t,s)$ be the unitary propagator
solving
\begin{equation}\label{eq:fiber-evolution}
 i\partial_t U_{\varepsilon,k}(t,s)
 =H_{k+\varepsilon t e_\beta}U_{\varepsilon,k}(t,s),
 \qquad U_{\varepsilon,k}(s,s)=I.
\end{equation}
The common domain is $\DD$. The temporal gauge gives the exact current
formula
\begin{equation}\label{eq:exact-current}
 j_\varepsilon(t)=\cB\int_\B
 \Tr\!\left[
 \bigl(v_\alpha(k)+\varepsilon t\,e_\alpha\cdot e_\beta\bigr)
 U_{\varepsilon,k}(t,0)P(k)U_{\varepsilon,k}(t,0)^*
 \right]dk.
\end{equation}
In particular the accelerated momentum is $p+k+\varepsilon t e_\beta$.
Both the sign and the scalar term in \eqref{eq:exact-current} are
consequences of the scalar potential $-\varepsilon x\cdot e_\beta$.
\Cref{lem:gauge,lem:local-trace} prove the formula for the regularity of
$V$ assumed here.

Let $\mathcal C$ denote complex conjugation. Since $V$ is real,
\begin{equation}\label{eq:time-reversal}
 \mathcal C H_k\mathcal C=H_{-k},\qquad
 \mathcal C P(k)\mathcal C=P(-k),\qquad
 \mathcal C v_\eta(k)\mathcal C=-v_\eta(-k).
\end{equation}
Thus $\Tr(v_\alpha(k)P(k))$ is real and odd under $k\mapsto-k$ on the
torus. Formula \eqref{eq:exact-current} at $\varepsilon=0$ proves that
$j_0(t)=0$.

\begin{proposition}\label{prop:raw-response}
For each fixed $t\geq0$, the derivative at zero field exists and equals
\begin{equation}\label{eq:raw-response}
 \begin{split}
 \sigma_{\alpha\beta}(t)=\cB\int_\B\Bigg\{&
 Nt\,e_\alpha\cdot e_\beta\\
 &-2\sum_{n\leq N<m}\Ree z_{nm}
 \left(\frac{t}{\Delta_{mn}}
       -\frac{\sin(t\Delta_{mn})}{\Delta_{mn}^2}\right)\\
 &+2\sum_{n\leq N<m}\Imm z_{nm}
 \frac{1-\cos(t\Delta_{mn})}{\Delta_{mn}^2}
 \Bigg\}\,dk.
 \end{split}
\end{equation}
All displayed band sums and integrals converge absolutely for fixed
$t$; the same bounds hold uniformly on bounded time intervals.
\end{proposition}

\begin{proof}
The differentiation of \eqref{eq:exact-current}, including the
unbounded factor $v_\alpha$, is justified in \cref{lem:differentiation}.
It gives
\begin{equation}\label{eq:commutator-response}
 \sigma_{\alpha\beta}(t)=\cB\int_\B
 \left\{Nt\,e_\alpha\cdot e_\beta
 -i\int_0^t s\,
 \Tr\!\left(
 v_\alpha e^{-i(t-s)H_k}[v_\beta,P]e^{i(t-s)H_k}
 \right)ds\right\}dk.
\end{equation}
For $r=t-s$, the trace in this formula is
\begin{equation}\label{eq:commutator-bands}
 \sum_{n\leq N<m}
 \bigl(z_{nm}e^{-ir\Delta_{mn}}
       -\overline{z_{nm}}e^{ir\Delta_{mn}}\bigr).
\end{equation}
Its absolute convergence follows from \eqref{eq:sum-bound}. The two
elementary integrals
\[
 \int_0^t s\sin((t-s)\Delta)\,ds
 =\frac{t}{\Delta}-\frac{\sin(t\Delta)}{\Delta^2},\qquad
 \int_0^t s\cos((t-s)\Delta)\,ds
 =\frac{1-\cos(t\Delta)}{\Delta^2}
\]
give \eqref{eq:raw-response}. Finally,
\[
 \frac{|\sin(t\Delta)|}{\Delta^2}\leq\frac{t}{\Delta},
 \qquad
 \frac{1-\cos(t\Delta)}{\Delta^2}\leq\frac{t^2}{2}.
\]
Together with \eqref{eq:sum-bound} and
\eqref{eq:inverse-gap-integrable}, these inequalities justify the
separate terms and all interchanges of sums and integrals.
\end{proof}

\section{Cancellation of the secular and antisymmetric terms}
\label{sec:cancellation}

The first two lines of \eqref{eq:raw-response} contain terms proportional
to $t$. Their cancellation is an identity for the total occupied band
energy. At the crossings, that energy is not twice differentiable in
the classical sense; the following argument accounts for those points.

\begin{lemma}[Integrated band-energy identity]\label{lem:sum-rule}
One has
\begin{equation}\label{eq:sum-rule}
 N|\B|\,e_\alpha\cdot e_\beta
 =2\int_\B\sum_{n\leq N<m}
 \frac{\Ree z_{nm}(k)}{\Delta_{mn}(k)}\,dk.
\end{equation}
\end{lemma}

\begin{proof}
On the punctured torus define
\[
 E(k)=\sum_{n=1}^N\lambda_n(k)=\Tr(H_kP(k)).
\]
The occupied cluster is isolated there, so its Riesz projection and
$E$ are smooth, even at internal band degeneracies. Differentiation of
the spectral projection gives
\begin{equation}\label{eq:energy-derivatives}
 \partial_\beta E=\Tr(v_\beta P),\qquad
 \partial_\alpha\partial_\beta E
 =N e_\alpha\cdot e_\beta
 -2\sum_{n\leq N<m}\frac{\Ree z_{nm}}{\Delta_{mn}}.
\end{equation}
For completeness, $P(\partial_\beta P)P=0$ and
$(I-P)(\partial_\beta P)(I-P)=0$. Differentiating $[H_k,P(k)]=0$
therefore yields, for $n\leq N<m$,
\[
 \ip{u_m}{(\partial_\beta P)u_n}
 =-\frac{\ip{u_m}{v_\beta u_n}}{\Delta_{mn}}.
\]
This and its adjoint give the second identity in
\eqref{eq:energy-derivatives}, since
$\partial_\beta v_\alpha=e_\alpha\cdot e_\beta$.

The first identity and \eqref{eq:occupied-bound} show that
$\nabla_k E$ is bounded on the punctured torus, uniformly up to the
punctures. The second
identity and \eqref{eq:inverse-gap-integrable} show that its classical
Hessian is integrable. Delete disks of radius $\rho$ around the
crossings and integrate $\partial_\alpha(\partial_\beta E)$ over the
remaining torus. Periodicity cancels the outer cell faces. Each small
circle contributes at most $C\rho$, because $\partial_\beta E$ is
bounded. Letting $\rho\downarrow0$ gives
\[
 \int_\B\partial_\alpha\partial_\beta E\,dk=0.
\]
Substitution of \eqref{eq:energy-derivatives} proves
\eqref{eq:sum-rule}. In particular, no point-supported correction is
left at a conical crossing.
\end{proof}

Time reversal removes the last line of \eqref{eq:raw-response}. Indeed,
at $-k$ we may use the eigenvectors $\mathcal C u_n(k)$, so that
\begin{equation}\label{eq:tr-matrix}
 \Delta_{mn}(-k)=\Delta_{mn}(k),\qquad
 z_{nm}(-k)=\overline{z_{nm}(k)}.
\end{equation}
These identities are understood after summation over any degenerate
eigenspaces. Reciprocal-lattice covariance makes the integrands
well defined on $\T$, whether or not the chosen cell $\B$ is symmetric.
The imaginary part in \eqref{eq:raw-response} is integrable and odd,
and hence its integral vanishes. The individual matrix elements need
not be real.

Combining this observation with \cref{lem:sum-rule} yields the response
formula used in the rest of the proof:
\begin{equation}\label{eq:response}
 \boxed{\displaystyle
 \sigma_{\alpha\beta}(t)
 =\frac{2}{(2\pi)^2}\int_\B\sum_{n\leq N<m}
 \Ree z_{nm}(k)\,
 \frac{\sin(t\Delta_{mn}(k))}{\Delta_{mn}(k)^2}\,dk.}
\end{equation}
The estimates already proved imply that this is a continuous function
of $t$ on every bounded time interval.

\section{The current matrix elements at a conical crossing}
\label{sec:cone}

The hypotheses describe band energies, while \eqref{eq:response}
depends on eigenvectors through $z_{nm}$. We now derive the required
matrix elements from the energy assumptions.

\begin{lemma}[Local conical matrix element]\label{lem:cone}
Fix a crossing $k_i$, write $v=v_i$, and let
$\omega=q/|q|$. Then, uniformly in $\omega\in S^1$,
\begin{equation}\label{eq:cone-matrix}
 \Ree z_{N,N+1}(k_i+q)
 =v^2\left(e_\alpha\cdot e_\beta
 -(e_\alpha\cdot\omega)(e_\beta\cdot\omega)\right)+O(|q|).
\end{equation}
\end{lemma}

\begin{proof}
The two bands meeting at $k_i$ form an isolated spectral cluster in a
neighborhood of $q=0$. Let $Q(q)$ be its rank-two Riesz projection.
This projection is real analytic into $\mathcal B(\HH,\DD)$. Indeed,
for $z$ on a fixed contour enclosing the pair, put
$R_0(z)=(H_{k_i}-z)^{-1}$. Then
\[
 (H_{k_i+q}-z)^{-1}
 =R_0(z)\bigl[I+(H_{k_i+q}-H_{k_i})R_0(z)\bigr]^{-1}.
\]
The inverse on the right is a uniformly convergent Neumann series for
small complex $q$, since $H_{k_i+q}-H_{k_i}$ is polynomial from $\DD$
to $\HH$. Contour integration proves the claim. This is the
isolated-cluster construction of analytic perturbation theory
\cite[Section XII.2]{RS4}; it does not assert analyticity of the
individual conical eigenvalues.

Choose an isometry $W_0:\C^2\to\Ran Q(0)$. For small real $q$, the map
\begin{equation}\label{eq:analytic-frame}
 W(q)=Q(q)W_0\bigl(W_0^*Q(q)W_0\bigr)^{-1/2}
\end{equation}
is a real analytic orthonormal frame for $\Ran Q(q)$, with values in
$\DD$. Define the Hermitian matrix
\begin{equation}\label{eq:pauli-h}
 h(q)=W(q)^*H_{k_i+q}W(q)
 =a(q)I+d(q)\cdot\boldsymbol\sigma,
\end{equation}
where $a(q)\in\R$, $d(q)\in\R^3$, and
$\boldsymbol\sigma=(\sigma_1,\sigma_2,\sigma_3)$ are the Pauli matrices.
Explicitly,
\[
 \sigma_1=\begin{pmatrix}0&1\\1&0\end{pmatrix},\quad
 \sigma_2=\begin{pmatrix}0&-i\\i&0\end{pmatrix},\quad
 \sigma_3=\begin{pmatrix}1&0\\0&-1\end{pmatrix}.
\]
The eigenvalues of $h$ are $a\pm|d|$. Thus \eqref{eq:cones} gives
\begin{equation}\label{eq:pauli-expansion}
 a(q)=\mu_{\mathrm F}+O(|q|^2),\qquad
 d(q)=Lq+O(|q|^2),\qquad L^*L=v^2I_{\R^2},
\end{equation}
where $L:\R^2\to\R^3$ is linear. To see the last assertion, expand the
analytic function $d$ at zero and compare
$|d(r\omega)|=vr+O(r^2)$ along each ray. It follows that
$|L\omega|=v$ for every unit $\omega$, which is precisely the asserted
identity. Derivative estimates for the remainders in
\eqref{eq:pauli-expansion} follow from analyticity of $a$ and $d$.

The physical velocities compressed to this frame are
\[
 B_\eta(q)=W(q)^*v_\eta(k_i+q)W(q).
\]
They must be distinguished from the derivatives of $h$, because the
frame depends on $q$. With $A_\eta=W^*\partial_\eta W$, differentiation
of $H_{k_i+q}W=Wh$ and $W^*W=I$ gives
\begin{equation}\label{eq:connection}
 \partial_\eta h=B_\eta+[h,A_\eta].
\end{equation}
The matrix $A_\eta$ is bounded, and $h(0)=\mu_{\mathrm F}I$. Therefore
\begin{equation}\label{eq:velocity-expansion}
 B_\eta(q)=(Le_\eta)\cdot\boldsymbol\sigma+O(|q|).
\end{equation}

For $q\ne0$, set $n(q)=d(q)/|d(q)|$ and
$\pi_\pm(q)=(I\pm n(q)\cdot\boldsymbol\sigma)/2$. The active bands
are simple near the crossing, so
\begin{equation}\label{eq:matrix-trace}
 z_{N,N+1}(k_i+q)
 =\tr_{\C^2}\bigl(\pi_- B_\alpha\pi_+B_\beta\bigr).
\end{equation}
For real $a,b,n\in\R^3$ with $|n|=1$, Pauli multiplication gives
\begin{equation}\label{eq:pauli-identity}
 \Ree\tr\!\left[
 \frac{I-n\cdot\boldsymbol\sigma}{2}
 (a\cdot\boldsymbol\sigma)
 \frac{I+n\cdot\boldsymbol\sigma}{2}
 (b\cdot\boldsymbol\sigma)\right]
 =a\cdot b-(a\cdot n)(b\cdot n).
\end{equation}
Since $n(r\omega)=L\omega/v+O(r)$, insertion of
\eqref{eq:velocity-expansion} into \eqref{eq:matrix-trace} and use of
$L^*L=v^2I$ prove \eqref{eq:cone-matrix}.
\end{proof}

The angular factor in \eqref{eq:cone-matrix} will be denoted by
\begin{equation}\label{eq:angular-factor}
 f_{\alpha\beta}(\omega)
 =e_\alpha\cdot e_\beta
 -(e_\alpha\cdot\omega)(e_\beta\cdot\omega).
\end{equation}
With $d\omega$ denoting arc length on $S^1$,
\begin{equation}\label{eq:angular-integral}
 \int_{S^1}f_{\alpha\beta}(\omega)\,d\omega
 =\pi\,e_\alpha\cdot e_\beta.
\end{equation}
Indeed, $\int_{S^1}\omega\otimes\omega\,d\omega=\pi I_{\R^2}$.

\section{The Ces\`aro limit}\label{sec:limit}

Time integration of \eqref{eq:response} gives
\begin{equation}\label{eq:cesaro}
 \frac1T\int_0^T\sigma_{\alpha\beta}(t)\,dt
 =\frac{2}{(2\pi)^2}\int_\B\sum_{n\leq N<m}
 \Ree z_{nm}(k)\,
 \frac{1-\cos(T\Delta_{mn}(k))}{T\Delta_{mn}(k)^3}\,dk.
\end{equation}
Fubini is justified by \eqref{eq:inverse-gap-integrable}, since
\[
 \frac1T\int_0^T\frac{|\sin(t\Delta)|}{\Delta^2}\,dt
 \leq\frac{T}{2\Delta}.
\]

Keep the fixed disks $D_i$ from \cref{sec:response}. All transitions
outside those disks, and all transitions inside them except the pair
$(N,N+1)$, have a common positive gap. By \eqref{eq:sum-bound}, their
total contribution to \eqref{eq:cesaro} is $O(T^{-1})$. It remains to
evaluate, for one crossing and a fixed sufficiently small $\delta>0$,
\begin{equation}\label{eq:cone-integral}
 I_i(T)=\frac{2}{(2\pi)^2}\int_{|q|<\delta}
 b_i(q)\,
 \frac{1-\cos(T\Delta_i(q))}{T\Delta_i(q)^3}\,dq,
 \qquad b_i(q)=\Ree z_{N,N+1}(k_i+q).
\end{equation}

Write $q=r\omega$, make the exact change of variables $s=Tr$, and put
$D_T(s,\omega)=T\Delta_i((s/T)\omega)$. Then
\begin{equation}\label{eq:scaled-integral}
 I_i(T)=\cB\int_{S^1}\int_0^\infty
 \one_{\{s<\delta T\}}
 \frac{2b_i((s/T)\omega)\,s}{D_T(s,\omega)^3}
 \bigl(1-\cos D_T(s,\omega)\bigr)\,ds\,d\omega.
\end{equation}
The integrand is taken to be zero when $s\geq\delta T$. On its support,
\eqref{eq:small-gap} gives $cs\leq D_T(s,\omega)\leq Cs$, and $b_i$ is
bounded. Consequently its absolute value is at most
\begin{equation}\label{eq:majorant}
 C\min\{1,s^{-2}\},
\end{equation}
an integrable function of $s$ independent of $T$ and $\omega$.

For every fixed $s>0$ and $\omega\in S^1$,
\[
 D_T(s,\omega)\longrightarrow2v_i s,\qquad
 b_i((s/T)\omega)\longrightarrow v_i^2 f_{\alpha\beta}(\omega).
\]
Dominated convergence in \eqref{eq:scaled-integral} therefore yields
\begin{equation}\label{eq:cone-limit-integral}
 \lim_{T\to\infty}I_i(T)
 =\cB\int_{S^1}\frac{f_{\alpha\beta}(\omega)}{4v_i}\,d\omega
 \int_0^\infty\frac{1-\cos(2v_i s)}{s^2}\,ds.
\end{equation}
The scalar integral is $\pi v_i$. One self-contained evaluation starts
from $u^{-2}=\int_0^\infty y e^{-yu}\,dy$: by Tonelli,
\begin{equation}\label{eq:scalar-integral}
 \int_0^\infty\frac{1-\cos u}{u^2}\,du
 =\int_0^\infty y\left(\frac1y-\frac{y}{1+y^2}\right)dy
 =\frac\pi2.
\end{equation}
Using \eqref{eq:angular-integral} in \eqref{eq:cone-limit-integral}, we
obtain
\begin{equation}\label{eq:one-cone}
 \lim_{T\to\infty}I_i(T)
 =\frac1{(2\pi)^2}\,\frac\pi4\,
   \pi e_\alpha\cdot e_\beta
 =\frac1{16}\,e_\alpha\cdot e_\beta.
\end{equation}
Summing over the $M$ crossings and adding the vanishing contribution
from the other transitions proves \cref{thm:main}.

The calculation also explains the independence from $v_i$: the two
current matrix elements supply a factor $v_i^2$, while the radial gap
and its integration supply the inverse factor. Ces\`aro averaging is
what produces the integrable bound \eqref{eq:majorant}; no assumption
on critical points of the other bands is needed.

\appendix
\section{Operator and dynamical details}\label[appendix]{sec:appendix}

This appendix verifies the analytic facts used to pass from the Stark
evolution to the response formula. Constants may depend on a fixed
compact set of momenta and on bounded time and field intervals.

\subsection{Domains and self-adjointness}

\begin{lemma}\label{lem:operators}
The operators in \eqref{eq:hamiltonians} are self-adjoint on the stated
domains. The fibers have compact resolvent, depend analytically on $k$
as operators from $\DD$ to $\HH$, and satisfy, uniformly on compact
sets of $k$,
\begin{equation}\label{eq:graph-bound}
 C^{-1}\norm{u}_{H^2}
 \leq\norm{H_ku}+\norm{u}
 \leq C\norm{u}_{H^2}.
\end{equation}
Their quadratic-form domains are $H^1_{\mathrm{per}}(\Omega)$.
For every real $\varepsilon$, the restriction of $H_\varepsilon$ to
$C_c^\infty(\R^2)$ is essentially self-adjoint.
\end{lemma}

\begin{proof}
The periodic extension of $V$ is uniformly locally $L^2$. In dimension
two, multiplication by such a potential is infinitesimally
$-\Delta$-bounded \cite[Theorem XIII.96]{RS4}:
\begin{equation}\label{eq:v-relative}
 \norm{Vu}_2\leq\eta\norm{\Delta u}_2+C_\eta\norm{u}_2
 \qquad(\eta>0).
\end{equation}
The same bound holds on the periodic cell, by the localized estimate
or by Sobolev interpolation. Kato--Rellich
\cite[Theorem X.12]{RS2} gives self-adjointness and
\eqref{eq:graph-bound}; the first-order terms in $k$ are also
infinitesimally $-\Delta$-bounded. Compactness of
$H^2_{\mathrm{per}}(\Omega)\hookrightarrow L^2_{\mathrm{per}}(\Omega)$
gives compact resolvent. The polynomial expression for $H_k$ proves
analyticity on its fixed domain. Localized Sobolev interpolation also
gives
\[
 |\ip{u}{Vu}|\leq\eta\norm{\nabla u}_2^2+C_\eta\norm{u}_2^2,
\]
on the full space and on the cell. This identifies the form domains and
gives a uniform lower bound on compact sets of momenta.

For the Stark operator, apply the Faris--Lavine theorem
\cite[Theorem X.38]{RS2} to
\[
 2H_\varepsilon=-\Delta+U+W,\qquad
 U(x)=-2\varepsilon x\cdot e_\beta,\qquad W(x)=2V(x).
\]
Choose $c>0$ and $d$ so that $U\geq-c|x|^2-d$. The domain
$C_c^\infty(\R^2)$ is dense and invariant under multiplication by $x_j$
and differentiation. The comparison operator
$-\Delta+U+W+2c|x|^2$ is essentially self-adjoint on this domain by
\cite[Theorem X.29]{RS2}: after adding a constant,
$U+2c|x|^2$ is nonnegative and locally $L^2$, while $W$ has relative
$-\Delta$ bound zero. Moreover, $-a\Delta+W$ is bounded below for any
fixed $0<a<1$, by the form estimate above. These are the two hypotheses
of Theorem X.38, which proves the assertion. In particular, no global
$L^p(\R^2)$ integrability of the periodic potential is used.
\end{proof}

The graph estimate and the bounded occupied energies also give
\begin{equation}\label{eq:projector-h2}
 \mathop{\mathrm{ess\,sup}}_{k\in\B}
 \norm{P(k)}_{\mathcal B(\HH,\DD)}<\infty.
\end{equation}

\subsection{The temporal gauge}

\begin{lemma}\label{lem:gauge}
The propagators in \eqref{eq:fiber-evolution} exist, preserve $\DD$,
and satisfy
\begin{equation}\label{eq:propagator-graph}
 \sup_{\substack{k\in\B,\,0\leq s,t\leq t_*\\
                  |\varepsilon|\leq\varepsilon_*}}
 \norm{U_{\varepsilon,k}(t,s)}_{\mathcal B(\DD)}<\infty
\end{equation}
for every finite $t_*,\varepsilon_*$. Let
$U_\varepsilon(t,s)$ be their direct integral, transformed back to
$L^2(\R^2)$. Then
\begin{equation}\label{eq:gauge-identity}
 e^{-itH_\varepsilon}
 =G_\varepsilon(t)U_\varepsilon(t,0),\qquad
 G_\varepsilon(t)=e^{i\varepsilon t x\cdot e_\beta}.
\end{equation}
\end{lemma}

\begin{proof}
The Hamiltonians $K(t)=H_{k+\varepsilon t e_\beta}$ have common domain
$\DD$ and are continuously differentiable in
$\mathcal B(\DD,\HH)$. Their graph norms are uniformly equivalent by
\eqref{eq:graph-bound}, and
\[
 \partial_t K(t)=\varepsilon v_\beta(k+\varepsilon t e_\beta).
\]
Thus $\partial_tK(t)(K(s)-i)^{-1}$ is bounded and norm continuous,
uniformly on the compact parameter sets under consideration. The
common-domain evolution theorem \cite[Theorem X.70]{RS2}, applied to
$1+iK(t)$ and with the scalar damping removed, gives the propagator.
For initial data in $\DD$, self-adjointness of $K(t)$ implies
$\partial_t\norm{U_{\varepsilon,k}(t,s)f}^2=0$. Density and the
backward construction then give unitarity.
The graph estimate in the proof of that theorem gives
\eqref{eq:propagator-graph}; its constants are uniform because
\[
 \norm{(K(t)-K(s))(K(s)-i)^{-1}}\leq C|t-s|.
\]
The stepwise propagators used in that construction are measurable in
$k$ and respect reciprocal-lattice covariance; both properties pass to
their strong limit.
The corresponding full-space evolution preserves $H^2(\R^2)$ and
solves the equation with Hamiltonian
$\tfrac12(p+\varepsilon t e_\beta)^2+V$.

To identify this evolution with the Stark group, take $f\in H^2(\R^2)$
and put $w(t)=G_\varepsilon(t)U_\varepsilon(t,0)f$. For
$\phi\in C_c^\infty(\R^2)$, the vector
$G_\varepsilon(t)^*\phi$ is differentiable in $H^2$ and has compact
support. Differentiating the scalar product, using
$G_\varepsilon(t)^*pG_\varepsilon(t)=p+\varepsilon t e_\beta$, gives
\begin{equation}\label{eq:weak-stark}
 \frac{d}{dt}\ip{\phi}{w(t)}
 =-i\ip{H_\varepsilon\phi}{w(t)}.
\end{equation}
The integrated identity extends to every $\phi\in D(H_\varepsilon)$
because $C_c^\infty$ is a core. To see uniqueness, fix
$\phi\in D(H_\varepsilon^2)$ and set
$\phi_s=e^{-i(s-t)H_\varepsilon}\phi$. The extended weak identity gives
$\partial_s\ip{\phi_s}{w(s)}=0$, whence
\[
 \ip{\phi}{w(t)}=\ip{e^{itH_\varepsilon}\phi}{f}
 =\ip{\phi}{e^{-itH_\varepsilon}f}.
\]
Thus $w(t)=e^{-itH_\varepsilon}f$. Extending in the initial data by
continuity proves \eqref{eq:gauge-identity} on all of $L^2$.
This weak argument requires no position moment of $f$.
\end{proof}

\subsection{The local current trace}

\begin{lemma}\label{lem:local-trace}
For every finite $\varepsilon,t$, the operator
$J_\alpha\gamma_\varepsilon(t)$ has a bounded periodic extension and
is locally trace class. Its trace per unit area is
\eqref{eq:exact-current}.
\end{lemma}

\begin{proof}
We first record a consequence of the normalization
\eqref{eq:bloch-transform}. If $A$ is a bounded periodic operator with
Hilbert--Schmidt fibers and
$\int_\B\norm{A(k)}_{\mathfrak S_2}^2\,dk<\infty$, then
\begin{equation}\label{eq:hs-localization}
 \norm{A\one_\Omega}_{\mathfrak S_2}^2
 =\norm{\one_\Omega A}_{\mathfrak S_2}^2
 =\frac1{|\B|}\int_\B\norm{A(k)}_{\mathfrak S_2}^2\,dk.
\end{equation}
To prove the first equality with the integral, use an orthonormal
basis of $L^2(\Omega)$ extended by zero to the full space. Its Bloch
transform at $k$ is that basis multiplied by $e^{-ik\cdot x}$, which
is again an orthonormal basis of the fiber. Parseval and Tonelli give
the assertion. Applying the result to $A^*$ gives the other equality.

Let $Q$ be a periodic orthogonal projection whose fibers have rank $N$
and are uniformly bounded as maps from $\HH$ to
$H^1_{\mathrm{per}}(\Omega)$. Then $J_\alpha Q$ is
bounded, and its fibers have uniformly bounded Hilbert--Schmidt norms.
By \eqref{eq:hs-localization}, both $\one_\Omega J_\alpha Q$ and
$Q\one_\Omega$ are Hilbert--Schmidt. Their product
\[
 \one_\Omega J_\alpha Q\one_\Omega
 =(\one_\Omega J_\alpha Q)(Q\one_\Omega)
\]
is trace class \cite[Theorem VI.22(h)]{RS1}. Evaluating the trace with
the same cell basis gives
\begin{equation}\label{eq:current-trace-normalized}
 \frac1{|\Omega|}\Tr(\one_\Omega J_\alpha Q\one_\Omega)
 =\cB\int_\B\Tr(v_\alpha(k)Q(k))\,dk.
\end{equation}
The exchange of the sum and integral is legitimate because the fiber
trace norms are integrable. Localization to any bounded set follows
by covering it with finitely many lattice cells and using the same
Hilbert--Schmidt factorization.

Apply this to
$Q_\varepsilon(t)=U_\varepsilon(t,0)\gamma_0U_\varepsilon(t,0)^*$.
Its fibers have rank $N$ and are uniformly bounded from $\HH$ to $\DD$ by
\eqref{eq:projector-h2} and \eqref{eq:propagator-graph}.
Equation \eqref{eq:gauge-identity} gives
\[
 \gamma_\varepsilon(t)
 =G_\varepsilon(t)Q_\varepsilon(t)G_\varepsilon(t)^*.
\]
Conjugation by $G_\varepsilon(t)$ preserves periodicity: it commutes
with each lattice translation up to a scalar phase, which cancels in
this expression. It also preserves $H^2$ for fixed $t$. Since
$G_\varepsilon(t)$ commutes with $\one_\Omega$ and
\[
 G_\varepsilon(t)^*J_\alpha G_\varepsilon(t)
 =J_\alpha+\varepsilon t\,e_\alpha\cdot e_\beta,
\]
unitary conjugation of the localized trace and
\eqref{eq:current-trace-normalized} prove
\eqref{eq:exact-current} and the asserted local trace class.
\end{proof}

\Needspace{9\baselineskip}
\subsection{Differentiating the unbounded current}

\begin{lemma}\label{lem:differentiation}
For fixed $t$, one may differentiate \eqref{eq:exact-current} at
$\varepsilon=0$ under the integral. The result is
\eqref{eq:commutator-response}.
\end{lemma}

\begin{proof}
Fix $k$ outside the crossings and a unit vector $u\in\Ran P(k)$. Set
$\psi_\varepsilon(t)=U_{\varepsilon,k}(t,0)u$,
$\psi_0(t)=e^{-itH_k}u$, and
$\delta_\varepsilon(t)=\psi_\varepsilon(t)-\psi_0(t)$.
Duhamel's formula on the common domain gives
\begin{equation}\label{eq:duhamel}
 \delta_\varepsilon(t)
 =-i\int_0^t U_{\varepsilon,k}(t,s)
 \left(\varepsilon s v_\beta(k)
       +\frac{\varepsilon^2s^2}{2}|e_\beta|^2\right)
 e^{-isH_k}u\,ds.
\end{equation}
Since occupied unit vectors have uniformly bounded $H^2$ norms,
\begin{equation}\label{eq:delta-bounds}
 \norm{\delta_\varepsilon(t)}_2\leq C_t|\varepsilon|,
 \qquad
 \norm{\delta_\varepsilon(t)}_{H^2}\leq C_t
 \quad(|\varepsilon|\leq1),
\end{equation}
uniformly in $k$ and $u$. The $L^2$ estimate holds for arbitrary $H^2$
initial data with a factor $\norm{u}_{H^2}$ on the right. Density then
gives strong convergence of $U_{\varepsilon,k}(t,s)$ to
$e^{-i(t-s)H_k}$, uniformly for $s\in[0,t]$.
It follows from \eqref{eq:duhamel} that
\begin{equation}\label{eq:vector-derivative}
 \frac{\delta_\varepsilon(t)}{\varepsilon}
 \longrightarrow
 \eta(t):=-i\int_0^t s\,e^{-i(t-s)H_k}
 v_\beta(k)e^{-isH_k}u\,ds
 \quad\text{in }L^2.
\end{equation}

This $L^2$ convergence suffices for the current. Interpolation in
\eqref{eq:delta-bounds} gives
$\norm{\delta_\varepsilon(t)}_{H^1}\leq C_t|\varepsilon|^{1/2}$.
The exact expansion
\begin{equation}\label{eq:current-expansion}
 \begin{split}
 \ip{\psi_\varepsilon}{v_\alpha\psi_\varepsilon}
 -\ip{\psi_0}{v_\alpha\psi_0}
 &=2\Ree\ip{\delta_\varepsilon}{v_\alpha\psi_0}
   +\ip{\delta_\varepsilon}{v_\alpha\delta_\varepsilon}
 \end{split}
\end{equation}
has a last term bounded by $C_t|\varepsilon|^{3/2}$. Thus its derivative
is $2\Ree\ip{\eta(t)}{v_\alpha\psi_0(t)}$. The bounds are uniform in
$k$ and in occupied unit vectors, so the difference quotients of the
fiber current have a uniform majorant after summing the $N$ occupied
states. Dominated convergence justifies differentiation of
\eqref{eq:exact-current}.

Summing \eqref{eq:vector-derivative} over an occupied orthonormal basis
gives the density derivative
\[
 -i\int_0^t s\,e^{-i(t-s)H_k}
 [v_\beta(k),P(k)]e^{i(t-s)H_k}\,ds.
\]
The scalar term in \eqref{eq:exact-current} contributes
$Nt\,e_\alpha\cdot e_\beta$, giving
\eqref{eq:commutator-response}. The trace in that formula is well
defined: $[v_\beta,P]$ is a bounded finite-rank operator with range in
$H^1$, and $e^{-irH_k}$ preserves $H^1$, the form domain from
\cref{lem:operators}. Multiplication on the left by $v_\alpha$ therefore
still yields a finite-rank bounded operator. Its spectral expansion is
the absolutely convergent series \eqref{eq:commutator-bands}.
\end{proof}

\begin{thebibliography}{9}
\bibitem{RS1}
M.~Reed and B.~Simon,
\emph{Methods of Modern Mathematical Physics. I: Functional Analysis},
revised and enlarged edition, Academic Press, New York, 1980.

\bibitem{RS2}
M.~Reed and B.~Simon,
\emph{Methods of Modern Mathematical Physics. II: Fourier Analysis,
Self-Adjointness}, Academic Press, New York, 1975.

\bibitem{RS4}
M.~Reed and B.~Simon,
\emph{Methods of Modern Mathematical Physics. IV: Analysis of Operators},
Academic Press, New York, 1978.
\end{thebibliography}
\end{document}
